% Part: sets-functions-relations
% Chapter: arithmetization
% Section: cauchy

\documentclass[../../../include/open-logic-section]{subfiles}

\begin{document}

\olfileid[pt]{sfr}{arith}{cauchy}

\olsection{Apêndice: os Reais como Sequências de Cauchy}

Na \olref[cuts]{sec}, construímos os reais como cortes de Dedekind. Nesta
secção, explicamos uma construção alternativa. Apoia-se na definição, devida a Cauchy,
do que hoje chamamos \emph{sequência de Cauchy}; mas o uso desta definição para
\emph{construir} os reais deve-se a outros autores do século XIX, nomeadamente
Weierstrass, Heine, M\'{e}ray e Cantor. (Para uma boa história, ver
\citeauthor{OConnorRobertson:RN}
\citeyear{OConnorRobertson:RN}.)

Antes de chegarmos ao século XIX, vale a pena considerar Simon Stevin
(1548--1620). Em resumo, Stevin apercebeu-se de que podemos pensar cada real em
termos da sua expansão decimal. Assim, mesmo um número irracional, como
$\sqrt{2}$, tem uma expansão decimal agradável, que começa:
\[
	1.41421356237\ldots
\] 
É muito fácil modelar expansões decimais na teoria dos conjuntos: basta
considerá-las como funções $d \colon \Nat \to \Nat$, onde $d(n)$ é o $n$-ésimo
algarismo decimal que nos interessa. Depois precisamos de um pequeno ajuste,
para tratar a parte do número real que vem antes da vírgula decimal (aqui,
apenas $1$). Precisaremos também de outro ajuste (uma relação de equivalência)
para garantir que, por exemplo, $0.999\ldots
= 1$. Mas não é difícil apresentar uma construção perfeitamente rigorosa dos
números reais, à maneira de Stevin, dentro da teoria dos conjuntos.

Stevin não é o nosso foco. (Para mais sobre Stevin, ver
\citealt{KatzKatz2012}.) Mas eis um pensamento estreitamente relacionado. Em
vez de tratar diretamente a expansão decimal de $\sqrt{2}$, podemos antes
considerar uma \emph{sequência} de aproximações racionais cada vez mais precisas
a $\sqrt{2}$, olhando para as expansões cada vez mais precisas:
\[
1, 1.4, 1.414, 1.4142, 1.41421,\ldots
\]
A ideia de que os reais podem ser considerados via ``aproximações cada vez
melhores'' nos dá a base para outra sequência de insights, semelhante às
constatações que retivemos ao construir $\Rat$ a partir de $\Int$, ou $\Int$ a
partir de $\Nat$. Os insights básicos são estes:
\begin{enumerate}
	\item Todo o real pode escrever-se como uma expansão decimal (eventualmente
	infinita).
	\item A informação codificada por uma expansão decimal (eventualmente
	infinita) pode igualmente codificar-se por uma sequência de números
	racionais.
	\item Uma sequência de números racionais pode pensar-se como uma função de
	$\Nat$ para $\Rat$; basta definir $f(n)$ como o $n$-ésimo racional na
	sequência.
\end{enumerate}
Claro que nem \emph{qualquer} função de $\Nat$ para $\Rat$ nos dará um número
real. Por exemplo, considere-se esta função:
\[
	f(n) =\begin{cases}
			1 & \text{se }n\text{ é ímpar}\\
			0 &\text{se }n\text{ é par}
		\end{cases}
\]
Essencialmente, a preocupação é que a sequência $0,1,0,1,0,1,0,\ldots$
não parece ``convergir'' para nenhum real. Logo: para garantir que consideramos
sequências que de fato convergem para algum real, precisamos de restringir a
atenção a sequências que tenham algum \emph{limite}.

Já encontramos a ideia de limite em
\olref[his][set][limits]{sec}. Mas não podemos usar \emph{exatamente} a mesma
definição que ali. A expressão ``$(\forall \epsilon>0)$''
envolvia tacitamente quantificação sobre os números reais; e estávamos a
considerar limites de funções nos reais; logo, invocar essa definição seria
valer-nos dos números reais --- e eram precisamente eles que estávamos a
tentar \emph{construir}. Felizmente, podemos trabalhar com uma ideia estreitamente
relacionada com a de limite.
\begin{defn}\ollabel{def:CauchySequence}
  Uma função $f: \Nat \to \Rat$ é uma \emph{sequência de Cauchy} sse, para
  qualquer $\epsilon \in \Rat$ positivo, se tem $(\exists \ell \in
  \Nat)(\forall m, n > \ell)|f(m) - f(n)| < \epsilon$.
\end{defn}

A ideia geral de limite é a mesma de antes: se se quer um certo nível de
precisão (medido por~$\epsilon$), há uma ``região'' onde olhar (qualquer entrada
maior que~$\ell$). E é fácil ver que a nossa sequência $1$, $1.4$, $1.414$,
$1.4142$, $1.41421$\ldots tem limite: se se quer aproximar $\sqrt{2}$ com erro
inferior a $\nicefrac{1}{10^{n}}$, basta olhar para qualquer entrada depois da
$n$-ésima.

O pensamento óbvio seria então dizer que um número real \emph{é} simplesmente
qualquer sequência de Cauchy. Mas, como nas construções de $\Int$ e $\Rat$, isso
seria demasiado ingênuo: para um dado número real, várias sequências de Cauchy
diferentes ``indicam'' esse número. Uma forma simples de ver isto é a seguinte.
Dada uma sequência de Cauchy~$f$, defina-se $g$ como exatamente a mesma
função que~$f$, exceto que $g(0)\neq f(0)$. Como as duas sequências coincidem
em todo o lado depois do primeiro número, quereremos, em última análise, dizer
que têm o mesmo limite, no sentido de \olref{def:CauchySequence},
e portanto devem ser pensadas como ``definindo'' o mesmo real. Logo, devemos
pensar nestas sequências de Cauchy como o mesmo número real.

Consequentemente, voltamos a precisar de definir uma relação de equivalência sobre
as sequências de Cauchy, e identificar os números reais com \emph{classes} de
equivalência. Precisamos primeiro da ideia de uma função que tende para $0$ no
limite. Para qualquer função $h : \Nat \to \Rat$, digamos que \emph{$h$ tende
para $0$} sse, para qualquer $\epsilon \in \Rat$ positivo, se tem que
$(\exists \ell \in \Nat)(\forall n > \ell)|h(n)| <
\epsilon$.\footnote{Compare-se com a definição de $\lim_{x
\mathord{\rightarrow}\infty}f(x) = 0$ em
\olref[his][set][limits]{sec}.} Além disso, sendo $f$ e $g$ funções $\Nat \to \Rat$, definamos $(f-g)(n) = f(n) - g(n)$. Definimos agora:
\[
	f \Realequiv g \text{ sse $(f-g)$ tende para $0$}.
\]
Há que verificar que $\Realequiv$ é uma relação de equivalência; e de fato é.
Podemos então, se quisermos, definir os reais como as classes de equivalência,
sob $\Realequiv$, de todas as sequências de Cauchy de $\Nat \to \Rat$.

\begin{prob}
Seja $f(n) = 0$ para todo $n$. Seja $g(n) = \frac{1}{(n+1)^2}$. Mostre
que ambas são sequências de Cauchy e que o limite de ambas as funções é $0$, de
modo que também $f \sim_\Real g$. 
\end{prob}

Feito isto, como de costume escrevemos $\equivrep{f}{\Realequiv}$
para a classe de equivalência que tem $f$ como !!a{element}. Porém, para manter a
legibilidade, no que se segue omitimos o subscrito e escrevemos apenas
$\equivrep{f}{}$. Estipulamos também que, para cada $q \in \Rat$,
se tem $q_{\Real} = \equivrep{c_{q}}{}$, onde $c_{q}$ é a função constante
$c_q(n) = q$ para todo $n \in \Nat$. Definimos depois relações e operações
básicas sobre os reais, por exemplo:
\begin{align*}
	\equivrep{f}{} + \equivrep{g}{} &= 	\equivrep{(f + g)}{} \\
	\equivrep{f}{} \times 	\equivrep{g}{} &= \equivrep{(f \times g)}{} 
\end{align*}
onde $(f + g)(n) = f(n) + g(n)$ e $(f \times g)(n) = f(n) \times
g(n)$. Claro que também há que verificar que cada uma de $(f + g)$,
$(f-g)$ e $(f\times g)$ é sequência de Cauchy quando $f$ e $g$ o são;
mas o são, e deixamo-lo ao leitor.

Por fim, definimos uma noção de ordem. Dizemos que $\equivrep{f}{}$ é
\emph{positivo} sse tanto $\equivrep{f}{}\neq 0_\Rat$ como $(\exists
\ell \in \Nat)(\forall n > \ell)0 < f(n)$. Depois dizemos $\equivrep{f}{} <
\equivrep{g}{}$ sse $\equivrep{(g - f)}{}$ é positivo. Há que verificar que isto
está bem definido (isto é, que não depende da escolha da função ``representante''
na classe de equivalência). 
%\begin{align*}
%	\equivrep{f}{} \leq \equivrep{g}{} \text{ iff }&\equivrep{f}{} = \equivrep{g}{} \text{ or }\\
%	&(\exists \ell \in \Nat)(\forall n \in \Nat)(\ell < n \lif f(n) \leq g(n))
%\end{align*}
Mas tendo feito isto, é bastante fácil mostrar que obtemos as propriedades algébricas certas; isto é:
\begin{thm}\ollabel{thm:cauchyorderedfield}
As sequências de Cauchy constituem um corpo ordenado.
\end{thm}

\begin{proof}
Exercício. 
\end{proof}

\begin{prob}
Mostre que as sequências de Cauchy constituem um corpo ordenado.
\end{prob}

É mais difícil provar que os reais, assim construídos, têm a propriedade de
completude, pelo que daremos a prova.

\begin{thm}
Todo o conjunto não vazio de sequências de Cauchy com majorante tem supremo
(menor majorante).
\end{thm}

\begin{proof}[Esboço de prova] Seja $S$ um conjunto não vazio de sequências de
Cauchy com majorante. Então existe algum $p \in \Rat$ tal que
$p_{\Real}$ é majorante de $S$. Seja $r \in S$; então existe algum
$q \in \Rat$ tal que $q_{\Real} < r$. Logo, se existe supremo de $S$, este está
entre $q_\Real$ e $p_\Real$ (inclusive). 

Nos aproximaremos do supremo, abordando-o simultaneamente por baixo e por cima. Em
particular, definimos duas funções, $f, g \colon \Nat \to \Rat$, com a ideia de
que $f$ se irá aproximar do supremo por cima e $g$ por baixo. Começamos por definir:
\begin{align*}
	f(0) &= p \\
	g(0) &= q
\end{align*}
Então, sendo $a_n = \frac{f(n) + g(n)}{2}$, definamos:\footnote{Isto é uma
definição recursiva. Mas ainda \emph{não} demos qualquer razão para crer que
definições recursivas são legítimas.}
\begin{align*}
	f(n+1) &=
	\begin{cases}
		a_n &\text{se }(\forall h \in S)\equivrep{h}{} \leq (a_n)_\Real\\
		f(n)&\text{caso contrário}
	\end{cases}\\
	g(n+1) &=
	\begin{cases}
		a_n &\text{se }(\exists h \in S)\equivrep{h}{} \geq (a_n)_\Real\\
	 	g(n) &\text{caso contrário}
	\end{cases}
\end{align*}
Tanto $f$ como $g$ são sequências de Cauchy. (Isto verifica-se com relativa
facilidade, mas deixamo-lo como exercício.) Note-se que a função $(f-g)$ tende
para $0$, pois a diferença entre $f$ e $g$ reduz-se à metade a cada passo. Logo
$\equivrep{f}{} = \equivrep{g}{}$. 

Mostramos primeiro que $\equivrep{f}{}$ é majorante de $S$, isto é, que $(\forall h \in S)\equivrep{h}{} \leq \equivrep{f}{}$.  
(Invocaremos \olref{thm:cauchyorderedfield} ao longo do argumento.) Seja $h \in S$ e
suponhamos, por \emph{reductio}, que $\equivrep{f}{} < \equivrep{h}{}$, de modo que
$0_\Real < \equivrep{(h-f)}{}$. Como $f$ é uma sequência de Cauchy monotonamente
decrescente, existe algum $n \in \Nat$ tal que
$\equivrep{(c_{f(n)} - f)}{} < \equivrep{(h-f)}{}$. Logo:
\[
	(f(n))_\Real = \equivrep{c_{f(n)}}{} < \equivrep{f}{} + \equivrep{(h-f)}{} = \equivrep{h}{},
\]
o que contradiz o fato de que, por construção, $\equivrep{h}{} \leq (f(n))_\Real$.

Mostramos a seguir que $\equivrep{f}{} = \equivrep{g}{}$ é o supremo de $S$. Seja
$j$ qualquer sequência de Cauchy e suponhamos $\equivrep{j}{} < \equivrep{g}{}$.
Raciocinando como acima (usando o fato de $g$ ser \emph{crescente}), existe
$n \in \Nat$ tal que $\equivrep{j}{} < (g(n))_\Real$. Mas, por construção, existe
$h \in S$ tal que $(g(n))_\Real \leq \equivrep{h}{}$, logo $\equivrep{j}{} <
\equivrep{h}{}$ e portanto $\equivrep{j}{}$ não é majorante de $S$.  
\end{proof}

\end{document}
